\scanpage{157}
\mspage{486}
% position: upper continuation paragraph
très fort auprès des femmes dont j'ai été amant par la suite, en réveillant en
l'amante des pulsions masculines, qui durant toute sa vie avaient été soigneu-\linebreak
semant refoulées, et n'avaient trouvé à s'exprimer jusque là qu'``à la sauvette''
comme des sortes de bavures, indignes de figurer dans le vécu amoureux conscient

% position: upper long paragraph
Le vécu amoureux inconscient est riche de pulsions archétypes, dont une \unsure{de|des}
plus puissantes est celle du retour à la Mère, du retour dans le giron originel.
Un tel archétype est présent dans les couches profondes de l'expérience amoureu-\linebreak
se, chez l'homme et aussi chez la femme. En la femme, les résistances à la sa-\linebreak
tisfaction d'une telle pulsion dans le vécu amoureux du couple sont plus fortes
encore que chez l'homme, où elle se heurte à un tabou-clef, et non à deux comme
chez elle. Chez l'un comme chez l'autre, la satisfaction de ces pulsions \unsure{-}dans
le vécu commun reste souvent plus ou moins symbolique et surtout,
dérobé à la conscience. Quand un tel archétype et ce vécu remontent des couches
profondes jusqu'à la lumière du jour, dans le champ du regard conscient, ce
vécu aussitôt se transforme, il acquiert une dimension nouvelle. En même temps
se libèrent des énergies considérables, précédemment comprimées par les mécanis-\linebreak
mes répressifs, ou liées par les tâches de répression. L'effet est celui d'une
\mbox{l i b é r a t i o n} immédiate de la pulsion érotique, se manifestant par une
intensité renouvelée et par une plénitude nouvelle dans l'expérience amoureuse.

% position: lower long paragraph
Par ce qui précède, il apparaîtra déjà, sûrement, que cette acceptation
nouvelle de ma propre personne est allée de pair avec une acceptation d'autrui.
L'un et l'autre sont indissolublement liés. Il est entendu qu'il s'agit ici
d' ``acceptation'' dans le plein sens du terme, qui ne signifie nullement une
\mbox{t o l é r a n c e} (souvent aigre-douce) vis-à-vis de tels et tels ``travers'' ou
``défauts'', ressentis comme un mal hélas inévitable, pour lequel on est bien
obligé de ``faire avec''. Dans une telle attitude, je sens surtout une résignation,
pour ne pas dire une abdication, et sûrement pas une source de joie, ni un élan
de prise de connaissance d'une chose digne d'être connue: la profondeur préssen-\linebreak
tie\inserted{,}inconnue, derrière la plate surface de tels ``défauts'' ou ``travers'' qu'on
veut bien tolérer...

% position: bottom paragraph
Qu'il s'agisse ici d'une acceptation joyeuse, créatrice, ne signifie nulle-\linebreak
ment pour autant que cette acceptation soit totale-hier déjà j'ai constaté qu'il
n'en était rien. Un lecteur attentif l'aura déjà constaté par lui-même plus d'une
fois au cours de Récoltes et Semailles, comme il m'est arrivé de m'en
% position: left foot short mark
\margin{\unsure{-}}
% position: foot pencil
\margin{156}
